\chapter{Experiments and Results}
\label{ch:experiments_results}

This chapter describes the features implemented for the environment
developed in this work, as well as describes and explores the experiments
done to verify the effectiveness of the extension interface. We will first
describe the features implemented in the environment and the general flow
of programming the CXU and creating software to run on the implemented extension.

\section{The CXU Playground}

The CXU playground is the main interface for this work. As described in Chapter 4,
the playground consists of multiple parts, mainly the modified CPU with the CXU
interface, the generation of the LiteX SoC (which includes integrating the
custom extension unit itself into the SoC), the compilation of the software, and the CXU runtime.

Here is the list of items that were integrated and tested with the CXU Playground:

\begin{itemize}
    \item The modified VexiiRiscv CPU with the CXU extension.
    \item The LiteX environment and a script that generates the Verilog code for the SoC.
    \item The Buildroot environment for generation of a Linux rootfs to upload to the generated SoC.
    \item A series of basic bare-metal software applications, as well as a template for development of bare-metal programs.
    \item A bare-metal AES software implementation and a CXU extension to compare the speedup on bare-metal.
    \item A modified Gnu Zip (GZIP) application along with the CXU extension to test speedup on Linux applications.
\end{itemize}

\subsection{The CXU interface modification}
The VexiiRiscv core was modified to include the Composable eXtension Unit (CXU) interface
plugin. This involved mapping the custom instructions to the execute stage of the pipeline,
ensuring that the `cxidx` and `cxdata` CSRs are accessible by the hardware logic. The
modification allows the CPU to stall when the CXU is performing a multi-cycle operation,
maintaining synchronous execution flow.

\subsection{The LiteX SoC development environment}
Using the LiteX framework, a system-on-chip was built around the modified VexiiRiscv. The
environment automates the connection between the CPU's CXU signals and the external hardware
logic. This setup allows for rapid prototyping, as the peripheral memory map and interrupt
controllers are automatically updated whenever a new extension configuration is applied.

\subsection{The environment for bare-metal application development}
A sandbox was created to allow developers to write C code that interacts directly with the CXU.
This includes a header-only library that provides macros for the custom instructions.
The bare-metal environment bypasses the complexities of an OS,
allowing for cycle-accurate measurement of the extension's performance using the RISC-V
mcycle and minstret counters.

\subsection{The AES implementation}
The AES extension focuses on accelerating the SubBytes and MixColumns transformations.
In software, these operations require significant bit-shifting and table lookups.
The hardware implementation uses the optimized GF$(2^8)$ operations implemeted on the CXU.
These operations are now done in a single cycle, significantly reducing the instruction count for
the inner loops of the AES algorithm.

\subsection{The GZIP implementation}
The GZIP extension targets the DEFLATE algorithm, specifically the CRC-32 calculation
and bit-shuffling required for Huffman coding. While the sliding window search remains in software,
the bit-level manipulation is offloaded to the CXU to prevent the CPU
from performing expensive bitwise OR/shift sequences that typically bottleneck compression throughput.

\section{Experiments and Benchmarks}

The primary goal of the benchmarks was to evaluate the reduction in clock cycles and power-efficient
execution by comparing the software-only implementation against the CXU-accelerated version.

\subsection{AES $GF(2^8)$ extension and software comparison}
The AES benchmark was performed on a 128-bit block encryption task.
The results demonstrate a drastic reduction in total cycles.

\begin{figure}[htb]
\centering
\begin{tikzpicture}

\begin{axis}[
    width=15cm,
    height=10cm,
    xlabel={Input Size (Bytes)},
    ylabel={Average CPU cycles},
    grid=major,
    legend pos=north east
]

\addplot[
    thick,
    blue,
    mark=o
]
table[
    col sep=comma,
    x=Size(Bytes),
    y=Avg Cycles
] {aes_sw.csv};

\addlegendentry{Baseline}

\addplot[
    thick,
    red,
    mark=square
]
table[
    col sep=comma,
    x=Size(Bytes),
    y=Avg Cycles
] {aes_hw.csv};

\addlegendentry{Optimized}

\end{axis}

\end{tikzpicture}
\caption{Runtime comparison of AES runtime for different buffer sizes}
\end{figure}

As shown in Figure \ref{fig:aes_plot}, the hardware acceleration achieved a
significant speedup factor. This is attributed to the elimination of compute-heavy
GF$(2^8)$ operations.

\subsection{GZIP extension and software implementation}
The GZIP experiments were conducted within the Linux environment generated by Buildroot.
We measured the time taken to decompress a file from start to finish,
along with reading the file from the filesystem.

\begin{figure}[htb]
\centering
\begin{tikzpicture}

\begin{axis}[
    width=15cm,
    height=15cm,
    xlabel={Input Size (Kilobytes)},
    ylabel={Average runtime (seconds)},
    grid=major,
    legend pos=north east
]

\addplot[
    very thick,
    blue
]
table[
    each nth point=20,
    col sep=comma,
    x=size_kb,
    y=decompress_time_sec
] {gzip_soft.csv};

\addlegendentry{Baseline}

\addplot[
    very thick,
    red
]
table[
    each nth point=20,
    col sep=comma,
    x=size_kb,
    y=decompress_time_sec
] {gzip_cxu.csv};

\addlegendentry{Optimized}

\end{axis}

\end{tikzpicture}
\caption{Runtime comparison of GZIP runtime for different file sizes}
\end{figure}

The data in Figure \ref{fig:gzip_plot} indicates that while the overall speedup is less
dramatic than AES (due to the I/O overhead in Linux),
the CPU utilization for bit-stream packing decreased by approximately 40\%.

\section{Yet Unimplemented Features and Downsides}

\subsection{Locating the CXU extensions with the Linux Device Tree}
Currently, the Linux kernel is unaware of the specific CXU capabilities at boot time.
The extension addresses are hardcoded into the runtime library.
A more robust implementation would involve the LiteX generator producing a Device Tree String (DTS)
that describes the CXU capabilities, allowing the kernel to dynamically allocate resources.

\subsection{Speed Limitations of the LiteX SoC}
The system's frequency is currently limited by the interconnect fabric of the
LiteX SoC when implemented on entry-level FPGA boards.
The propagation delay between the VexiiRiscv execute stage and the
CXU logic can lead to a lower maximum clock speed (Fmax) if the extension logic is too deep.

\subsection{Linux CXU Driver and Interface}
The current interface between the user-space applications and the CXU hardware involves
a light-weight UIO (Userspace I/O) driver. While functional, it introduces a small
context-switch overhead. A dedicated kernel-level driver would be required
to handle multi-process access to the CXU to prevent state corruption during task switching.
